\begin{abstract}
A conformal prediction region promises to contain a road user's true future path a fixed fraction of the time. We ask
whether that promise survives when a trajectory forecaster is calibrated on one driving dataset and deployed, without
new labels, on another. Using forecasters calibrated on Argoverse 2 or nuScenes and evaluated on each other and on
Waymo Open Motion, we measure coverage losses of 6 to 25 percentage points at a 90 percent nominal level, on two
architectures whose accuracy is within 15 percent of published values; the third dataset was registered in advance.
Reweighting by an estimated covariate-shift likelihood ratio makes the loss worse, and we show that no label-free method
can repair an unrestricted shift, although the drift of the forecaster's own predicted scale tracks the loss. The loss
concentrates in turns, most of all in left-hand traffic. Measured input
differences between the datasets explain about a quarter of it for one architecture and, with caveats, all of it for
the other. Labelled target data can repair it, but labels arrive one scene at a time, and a guarantee that treats agents
as independent can fail when scenes contain many agents. A scene-level Learn-then-Test rule with a betting test restores
the guarantee at every labelling budget we tested, and returns a region within twice the oracle size from about 150 to
200 labelled scenes. Conformal guarantees for trajectory prediction should therefore not be assumed to transfer across
datasets without an audit and a scene-level recalibration.
\end{abstract}
